\documentclass[11pt]{article}

\usepackage[a4paper,margin=1in]{geometry}
\usepackage{amsmath,amssymb,amsthm,mathtools}
\usepackage{enumitem}
\usepackage{microtype}
\usepackage[hidelinks]{hyperref}

\newtheorem{theorem}{Theorem}[section]
\newtheorem{proposition}[theorem]{Proposition}
\newtheorem{lemma}[theorem]{Lemma}
\newtheorem{corollary}[theorem]{Corollary}
\newtheorem{definition}[theorem]{Definition}
\newtheorem{remark}[theorem]{Remark}

\newcommand{\Pref}{\operatorname{Pref}}
\newcommand{\rk}{\operatorname{rk}}
\newcommand{\RegAtlas}{\operatorname{RegAtlas}}
\newcommand{\Birth}{\operatorname{Birth}}
\newcommand{\Prov}{\operatorname{Prov}}
\newcommand{\Mother}{\operatorname{Mother}}
\newcommand{\supp}{\operatorname{supp}}
\newcommand{\PP}{\mathbb P}
\newcommand{\NN}{\mathbb N}

\title{The Prime-Distribution Birth Mother\\
\large Single-Prime, Finite-Combination, and All-Prime Return Births}
\author{}
\date{}

\begin{document}

\maketitle

\begin{abstract}
We begin with the unlabelled atomic successor history and the regular-atlas
mechanism extracted from the source manuscript.  Its one-direction primitive
births have posterior ranks exactly equal to the primes.

We prove first that the one-direction source alone cannot support a genuine
higher birth: its atlas refinements and intermediate covers recover exactly a
proper-divisor poset.  We then audit the tempting free cubical enlargement.
Although its boxes have genuine higher faces, its regular-grid atlas predicate
factors through a product of divisor posets.  The cubical face category and
relative fundamental class are identical for born and inherited boxes and do
not participate in the birth decision.  Thus the cubical proposal is also a
divisor shadow with an inert topological shell, not a joint prime birth.

Both failures are retained as provenance obstructions.  No probability law,
fixed prime spacing, batch operation, factorization tower, or coherence layer
is introduced.  A final finite-packet trial uses overlapping returns in one
atomic path trace.  It realizes every finite set of distinct primes, has
connected-support descent to the original unary births, and has mixed fibres
not determined by divisor data.  The universal unlabelled path-history source
then generates the return packets without receiving a prime list.  Its
simple-return sector has posterior support exactly all nonempty finite sets of
distinct primes.  On the inverse-limit boundary, an explicit
generated-subsource inheritance principle forces birth to mean dense orbit
under the canonical path shift.  Every source-transitive history contains
every finite joint birth as a window, and its primitive normalized-return
spectrum is the entire prime set.  The carrier histories are nonunique, but
they form one uniquely determined source-provenance class.
\end{abstract}

\section{Scope and provenance reset}

This paper studies only the birth mother of the prime distribution.  Its first
positive construction is the finite ordered-path source and the original
regular-atlas single-birth rule.  Endpoint covers, intermediate-cover nerves,
and cubical products appear only in obstruction audits; none is promoted to a
valid higher-birth source.  The later overlapping-return construction is
promoted through a universal unlabelled path-history source to a finite-set
birth theorem and then to a single-path all-prime boundary birth.  The boundary
uses only the canonical shift of that path source; no stack,
representation-theoretic, mapping-class, coherence, or probabilistic layer is
used in those constructions.

The word \emph{provenance} will have a narrow causal meaning.  It does not
refer to every structure from which the present model might be generalized.
It records only the data that were already present when a stage was born and
the path witnesses by which that stage was classified.

There are three levels that must not be conflated:
\begin{enumerate}[label=(\roman*)]
\item the \emph{source level}, consisting of unlabelled atomic extensions;
\item the \emph{single-birth level}, determined by the chronological
regular-atlas test against earlier primitive prefixes;
\item the \emph{arithmetic shadow}, obtained by counting atomic extensions
after the source and birth data have been formed.
\end{enumerate}
The arithmetic shadow is used to state and prove the identification with the
ordinary primes.  It is not an input to the mother rule.

\section{The unlabelled mother source}

\begin{definition}[Atomic successor source]
An \emph{unlabelled atomic successor source} is a rooted directed system
$(\mathscr R,o,\sigma)$ such that:
\begin{enumerate}[label=(\roman*)]
\item $o$ is the root;
\item every non-root state has one immediate predecessor;
\item every retained state has one next state supplied by the source rule
$\sigma$;
\item every state lies at the end of a finite chain of atomic successor moves
from $o$;
\item the source carries no numerical labels or arithmetic predicates.
\end{enumerate}
We assume that prefixes of every finite atomic depth occur.
\end{definition}

The source therefore produces one persistent history
\[
o\longrightarrow x_1\longrightarrow x_2\longrightarrow\cdots,
\]
where the subscripts are notation in the exposition and are not source data.

\begin{definition}[Retained prefix]
For a retained state $x$, let $H(x)$ be the entire directed chain from $o$ to
$x$.  Write $\Pref(\mathscr R)$ for the ordered family of nonempty retained
prefixes, ordered by initial-segment inclusion.
\end{definition}

The next candidate is always the next prefix produced by $\sigma$.  The
mother does not first choose an integer and then construct an interval of that
length.

\begin{definition}[Posterior rank]
For $H\in\Pref(\mathscr R)$, define
\[
\rk(H)
:=
\#\{\text{atomic source edges in }H\}.
\]
The rank is a posterior invariant of a retained path.  It is not available to
the birth predicate.
\end{definition}

\begin{proposition}[Clock reconstruction]
The prefix order and the atomic successor relation determine a unique order
isomorphism
\[
\Pref(\mathscr R)\cup\{H(o)\}\cong\NN_0.
\]
It is the map $H\mapsto\rk(H)$.
\end{proposition}

\begin{proof}
Every prefix contains a unique finite chain from the root.  Atomic extension
increases its edge count by one, and exactly one prefix occurs at each finite
depth.  Hence edge count gives the required order isomorphism.  Any order
isomorphism preserving the root and atomic successor must give the same
values.
\end{proof}

\section{Regular atlas inheritance}

We assume that the atomic edges have one local path type.  Thus two finite
segments are isomorphic precisely when their endpoints, order, and atomic
edge decompositions agree.

\begin{definition}[Regular atlas]
Let $\Pi$ be an earlier retained prefix and let $H$ be the current prefix.  A
\emph{regular $\Pi$-atlas} on $H$ is an ordered family
\[
(K_1,\ldots,K_r),\qquad r\ge 2,
\]
of consecutive segments of $H$ such that:
\begin{enumerate}[label=(\roman*)]
\item every $K_j$ is isomorphic to $\Pi$ as an ordered atomic path;
\item adjacent blocks meet only at their common endpoint;
\item the block interiors are pairwise disjoint;
\item the ordered concatenation of the blocks is all of $H$.
\end{enumerate}
Write $\RegAtlas(H;\Pi)$ for the finite set of these atlases.
\end{definition}

One occurrence of $\Pi$ inside $H$ is not an inheritance certificate.  The
copies must exhaust the current path.  On a linear homogeneous source, a
regular atlas of a fixed block type is unique whenever it exists.

\begin{lemma}[Rank of an atlas]
\label{lem:atlas-rank}
If $H$ has a regular $\Pi$-atlas with $r$ blocks, then
\[
\rk(H)=r\,\rk(\Pi).
\]
Conversely, if this equality holds for an integer $r\ge2$, then the unique
linear cuts of $H$ into $r$ consecutive blocks give a regular $\Pi$-atlas.
\end{lemma}

\begin{proof}
Each atlas block contains the same number of atomic source edges as $\Pi$.
The interiors do not overlap and the blocks exhaust $H$, so their edge counts
add.  Conversely, the displayed equality specifies the cuts after every
$\rk(\Pi)$ atomic edges.  Homogeneity identifies each resulting block with
$\Pi$.
\end{proof}

\section{One-direction provenance and mother relation}

\begin{definition}[Birth ledger]
The \emph{birth ledger before $H$}, denoted $\mathcal B_{<H}$, is the ordered
family of earlier prefixes declared primitive by the same rule below.  The
one-edge prefix \(E\) is the atomic unit and is not an arithmetic birth.
\end{definition}

\begin{definition}[Mother relation]
For $\Pi\in\mathcal B_{<H}$, define
\[
\Mother(\Pi,H)
\quad\Longleftrightarrow\quad
\RegAtlas(H;\Pi)\ne\varnothing.
\]
When this holds, $\Pi$ is an \emph{ancestral mother} of $H$.  A stage may have
more than one ancestral mother.
\end{definition}

\begin{definition}[Stage provenance]
The provenance record of the current stage is
\[
\Prov(H)
=
\left(
H,\ \mathcal B_{<H},\
\bigl(\RegAtlas(H;\Pi)\bigr)_{\Pi\in\mathcal B_{<H}}
\right).
\]
It contains the complete retained prefix, the earlier birth ledger, and the
actual atlas result for every earlier primitive prefix.  If a regular atlas
exists, that atlas is a positive inheritance witness.  If none exists, the
finite family of empty atlas sets is the stage's primitive-birth certificate.
\end{definition}

This definition makes provenance chronological and auditable.  Later
arithmetic names may describe the record, but they are not inserted into it.
In particular, the record does not contain a precomputed prime table, a
divisor list, or a factorization.

\begin{definition}[One-direction mother sector]
The \emph{one-direction sector} of the prime-distribution birth mother is
the system
\[
\mathfrak M
=
(\mathscr R,o,\sigma,\Pref(\mathscr R),\Prov,\Mother,\Birth)
\]
with the recursive law
\[
\boxed{
\Birth(H)=1
\quad\Longleftrightarrow\quad
\Mother(\Pi,H)\text{ fails for every }\Pi\in\mathcal B_{<H}.
}
\]
If $\Birth(H)=0$, the stage is inherited.  If $\Birth(H)=1$, $H$ is appended
to the ledger and becomes available as a possible ancestral mother of later
stages.
\end{definition}

In this sector the term \emph{mother} refers to the entire
source-and-ledger mechanism.  An
individual prime birth has no earlier ancestral mother; it is born precisely
at the first uncovered stage.  Once born, it becomes a mother direction for
future inherited stages.

\section{Exact prime-birth theorem}

\begin{theorem}[The birth ranks are exactly the primes]
\label{thm:single}
For every retained prefix $H$ with $\rk(H)\ge2$,
\[
\boxed{
\Birth(H)=1
\quad\Longleftrightarrow\quad
\rk(H)\in\PP.
}
\]
\end{theorem}

\begin{proof}
Use strong induction on the posterior rank.  The induction is part of the
external verification of the rule; it is not part of the rule's input.

Suppose first that $\rk(H)$ is prime.  If $H$ were inherited, some earlier
birth $\Pi$ would give a regular atlas with $r\ge2$ blocks.  The atlas-rank
lemma would give
\[
\rk(H)=r\,\rk(\Pi),
\]
with both factors greater than one, contradicting primality.  Thus $H$ is a
new birth.

Now suppose that $\rk(H)$ is composite.  Choose a prime $p$ dividing this
rank.  Clock reconstruction supplies a unique earlier prefix $\Pi$ of rank
$p$.  By the induction hypothesis, $\Pi$ is already in the birth ledger.
Writing $\rk(H)=rp$ with $r\ge2$, the converse part of the atlas-rank lemma
gives a regular $\Pi$-atlas on $H$.  Hence $H$ is inherited.  This completes
the induction.
\end{proof}

The proof uses ordinary arithmetic only to identify the numerical shadow of
the already-defined path recursion.  Operationally, the mother merely tests
the finitely many atlas sets recorded in $\Prov(H)$.

\begin{corollary}[Relabelling invariance]
An isomorphism of rooted atomic successor sources preserves the provenance
records, the mother relation, and the primitive births.  The construction is
therefore independent of names assigned to source states.
\end{corollary}

\section{The resulting prime distribution}

For a primitive prefix $\Pi$, define its future mother cone by
\[
\mathcal K(\Pi)
:=
\{H\in\Pref(\mathscr R):\Mother(\Pi,H)\}.
\]
At a current stage $H$, the mother rule is equivalently
\[
\Birth(H)=1
\quad\Longleftrightarrow\quad
H\notin\bigcup_{\Pi\in\mathcal B_{<H}}\mathcal K(\Pi).
\]

\begin{proposition}[Path-sieve identity]
After applying posterior rank, the cone of a primitive birth of rank $p$ is
the set of later ranks $2p,3p,4p,\ldots$.  The union of all earlier mother
cones at rank $n$ contains the rank-$n$ prefix exactly when $n$ is composite.
\end{proposition}

\begin{proof}
The atlas-rank lemma says that membership in the cone is equivalent to
$\rk(H)=r p$ for an integer $r\ge2$.  The exact prime-birth theorem identifies
the primitive ranks $p$ with the primes.  Every composite has a prime divisor,
and no prime belongs to a cone generated by an earlier primitive rank.
\end{proof}

Thus the mother is a path-native form of the sieve of Eratosthenes: each new
birth opens a future cone, while the uncovered stages are exactly the next
births.

\begin{definition}[Birth measure and birth-counting function]
Let $H_m$ denote the unique retained prefix of posterior rank $m$.  Define
\[
\mu_{\mathfrak M}
:=
\sum_{m\ge2}\Birth(H_m)\,\delta_m
\]
and
\[
B_{\mathfrak M}(x)
:=
\mu_{\mathfrak M}([2,x])
=
\sum_{2\le m\le x}\Birth(H_m).
\]
\end{definition}

\begin{theorem}[Exact distribution identity]
The numerical shadow of the birth mother satisfies
\[
\boxed{
\mu_{\mathfrak M}=\sum_{p\in\PP}\delta_p
}
\]
and, for every real $x\ge2$,
\[
\boxed{
B_{\mathfrak M}(x)=\pi(x).
}
\]
Equivalently,
\[
\supp(\mu_{\mathfrak M})=\PP.
\]
\end{theorem}

\begin{proof}
The coefficient of $\delta_m$ in $\mu_{\mathfrak M}$ is $1$ exactly when the
rank-$m$ prefix is a primitive birth.  By the exact prime-birth theorem this
occurs exactly when $m$ is prime.  Summing through $x$ gives the counting
identity.
\end{proof}

If the births are ordered as
\[
P_1<P_2<P_3<\cdots
\]
in the prefix order, then their posterior ranks are
\[
\rk(P_k)=p_k.
\]
Therefore every statistic formed solely from the support of the prime set can
be read from the birth history.  For example, consecutive birth gaps have
numerical shadows $p_{k+1}-p_k$.  This is an exact transfer of data, not a new
estimate for those statistics.

\section{Why the one-direction source cannot produce a genuine higher birth}

The regular-atlas mechanism above proves the single-prime theorem, but it does
not contain independent higher directions.  This is a structural limitation,
not a missing choice of notation.

Write \(H_n\) for the unique retained prefix of posterior rank \(n\).  Let
\(\mathcal A_n\) be the category generated by all nontrivial regular atlases on
\(H_n\) and their refinements.  Atlas isomorphisms are identified, since the
source is unlabelled and atomically homogeneous.

\begin{proposition}[Atlas refinement is the divisor poset]
\label{prop:atlas-divisor}
There is an order isomorphism
\[
\boxed{
\mathcal A_n
\cong
\{d:1<d<n,\ d\mid n\}^{\mathrm{op}},
}
\]
where an atlas by \(H_d\) corresponds to \(d\), and the unique refinement
arrow \(d\to e\) exists exactly when \(e\mid d\).
\end{proposition}

\begin{proof}
By Lemma~\ref{lem:atlas-rank}, a regular \(H_d\)-atlas on \(H_n\) exists exactly
when \(n=r d\) for some \(r\ge2\).  Atomic homogeneity makes its cut set unique:
the cuts occur after \(d,2d,\ldots,(r-1)d\) edges.  Thus atlas objects are
exactly the proper divisors of \(n\).  An \(H_d\)-block admits a regular
\(H_e\)-subatlas exactly when \(d=s e\); composing these identical subatlases
is precisely atlas refinement.  Hence the refinement order is exactly the
opposite of the usual divisibility order.
\end{proof}

The chronological birth ledger sees only those atlas objects whose block is
an earlier primitive birth.  After the prime-birth theorem, these are exactly
the prime divisors of \(n\).  They form only a collection of unary inheritance
witnesses.  If one enlarges from the birth ledger to all prefix atlases, or
equivalently to all endpoint-cover intermediates, the result is precisely
\(\mathcal A_n\), not an additional higher source.

\begin{corollary}[The endpoint-cover construction adds no higher provenance]
\label{cor:return-divisor}
Let \(\overline H_n\to\overline E\) be obtained by identifying the endpoints of
\(H_n\).  Its proper connected intermediate-cover category is isomorphic to
\(\mathcal A_n\).  Therefore its nerve, realization, and every chain called a
descent tower are functorial constructions on the proper-divisor poset.
\end{corollary}

\begin{proof}
Cutting an intermediate based circle cover at its return points gives a unique
regular atlas, and endpoint closure reverses this operation.  Cover
factorization corresponds to atlas refinement.  Apply
Proposition~\ref{prop:atlas-divisor}.
\end{proof}

In particular, if \(n\) is square-free with \(k\) posterior prime factors, the
proper-divisor poset is the proper part of the Boolean lattice on those
factors.  Its order complex is the barycentric subdivision of
\(\partial\Delta^{k-1}\), hence a sphere \(S^{k-2}\).  That sphere is real
topology, but its entire incidence was reconstructed from factor layers of one
uniserial path.  It is not evidence that a joint event was already present at
the source level.

\begin{theorem}[Uniserial provenance obstruction]
\label{thm:uniserial-obstruction}
Starting only from the atomic successor source \(\mathscr R\), its prefixes,
regular atlases, atlas refinements, and canonical endpoint closures, one
cannot obtain a genuine \(k\)-birth for \(k\ge2\) whose provenance contains
source-native incomparable faces.  Every proposed higher incidence category
formed from those operations factors through \(\mathcal A_n\), and every
proposed descent tower is therefore a factor-refinement chain in the
proper-divisor poset of the posterior rank.
\end{theorem}

\begin{proof}
The prefix category of \(\mathscr R\) is thin and totally ordered: between two
prefixes there is at most one extension map, and any two comparable extension
chains agree.  Regular-atlas sets are likewise empty or singletons.  The only
new arrows created by closing under atlas composition are refinement arrows,
which Proposition~\ref{prop:atlas-divisor} identifies with divisor
containments.  Endpoint closure is equivalent to the same refinement
category by Corollary~\ref{cor:return-divisor}.  Taking nerves, realizations,
subcategories, or chains changes the presentation but creates no independent
source face.  Thus no \(k\)-ary source event precedes the factor category from
which the alleged higher birth and descent are read.
\end{proof}

This theorem invalidates the earlier proposal
\[
\left|N\mathcal I(H_n)\right|\cong S^{k-2}
\quad\Longrightarrow\quad
\text{\emph{\(H_n\) is a \(k\)-birth}}
\]
The left side classifies a square-free factor lattice; it does not supply the
missing joint provenance on the right.

\section{Audit of the multidirectional source already present in the source manuscript}

The source manuscript later introduces a \(d\)-directional continuation
complex with independent atomic germs and retained commuting squares.  This
is genuinely non-thin, so it passes the first necessary test that the
one-direction source fails.  Its established conclusions, however, are more
limited than the finite-combination birth required here:

\begin{enumerate}[label=(\roman*)]
\item a path-primitive finite continuation cover has simple deck continuation
group and therefore posterior degree \(p\);
\item for fixed \(p\), the primitive covers form a projective family of
different continuation directions;
\item general finite covers admit composition towers whose successive fibres
have posterior prime cardinalities.
\end{enumerate}

None of these statements makes a set of distinct primes into one birth.  The
projective compatibility complex uses one coefficient field
\(\mathbb F_p\), so it concerns several directions carrying the same posterior
prime.  The composition tower starts from a nonprimitive composite cover and
resolves it into primitive steps; it is a factorization theorem, not a birth
theorem.  Rebranding either object as an arbitrary finite-prime birth would
repeat the provenance error.

The failure is not special to the monogenic circle cover.

\begin{proposition}[Intermediate-cover towers are always subgroup towers]
\label{prop:cover-subgroup}
Let \(q:Y\to X\) be a connected finite regular cover with deck group \(D\).
The category of connected intermediate covers
\[
Y\longrightarrow Z\longrightarrow X
\]
is equivalent, up to reversal of arrows, to the subgroup poset of \(D\).
Consequently every descent obtained solely by passing to intermediate covers
is a subgroup-series factorization of \(D\), even when the source \(X\) has
several continuation directions.
\end{proposition}

\begin{proof}
Fixing a basepoint identifies \(q_*\pi_1(Y)\) as a normal finite-index subgroup
of \(\pi_1(X)\), with quotient \(D\).  Connected intermediate covers correspond
to subgroups between \(q_*\pi_1(Y)\) and \(\pi_1(X)\).  Passing to the quotient
identifies these with subgroups of \(D\); cover composition reverses subgroup
inclusion.  A chain of intermediate covers is therefore exactly a subgroup
chain.
\end{proof}

Thus replacing the circle by the multidirectional continuation complex removes
thinness from the ambient source but does not by itself repair joint
provenance.  If the proposed \(k\)-birth and all its descents are read only from
intermediate covers of one regular cover, Proposition~\ref{prop:cover-subgroup}
shows that the result is still a factorization tower.  The missing higher faces
must come from the extension source itself, not from the lattice of a deck
quotient.

\section{The cubical product obstruction}

The free cubical completion is a legitimate source enlargement, but legitimacy
of the source object is not enough.  The birth predicate must itself use
higher topology rather than a product of unary inheritance tests.

For a framed box write
\[
B_{\mathbf n}=H_{n_1}\times\cdots\times H_{n_k},
\qquad
B_{\mathbf d}=H_{d_1}\times\cdots\times H_{d_k}.
\]
Consider the proposed regular cubical atlas: a Cartesian grid of consecutive
subboxes isomorphic to \(B_{\mathbf d}\), with one fixed matching of
parallelism classes.

\begin{proposition}[Cubical atlases factor coordinatewise]
\label{prop:cubical-product-divisor}
One has
\[
\boxed{
\RegAtlas_k(B_{\mathbf n};B_{\mathbf d})\ne\varnothing
\quad\Longleftrightarrow\quad
\exists\pi\in S_k\ 
\forall i,\ d_i\mid n_{\pi(i)},
}
\]
with at least one strict quotient.  Consequently the entire regular-grid
inheritance category factors through
\[
\left(\prod_{i=1}^k\mathcal D(n_i)\right)\big/S_k,
\]
where \(\mathcal D(n_i)\) is the divisor poset of the \(i\)-th posterior side
rank.
\end{proposition}

\begin{proof}
Intersect a cubical grid with an edge line in one parallelism class.  The
intersection is a one-dimensional regular atlas, so its block length divides
the corresponding side length.  Doing this in every direction gives the
right-hand condition.  Conversely, the one-dimensional cut sets determined by
those divisibilities have a Cartesian product, which is exactly the proposed
cubical grid.
\end{proof}

The resulting recursive classification
\[
\Birth_k(B_{\mathbf n})=1
\quad\Longleftrightarrow\quad
n_1,\ldots,n_k\text{ are prime}
\]
is therefore not an independent joint theorem.  It says that \(\mathbf n\) is
a minimal element in a product divisibility order.  Equivalently, it is the
conjunction of \(k\) unary prime-birth outcomes after a direction matching.

\begin{proposition}[The added cubical topology is inert]
\label{prop:cubical-inert}
The outer face category of every \(k\)-box has the same \(3^k\) face states,
independent of its side ranks, and every such box satisfies
\[
H_k(B_{\mathbf n},\partial B_{\mathbf n};\mathbb Z)
\cong\mathbb Z.
\]
Hence neither the face category nor the relative fundamental class
distinguishes a proposed born box from an inherited box.  Since neither enters
the regular-grid atlas predicate of
Proposition~\ref{prop:cubical-product-divisor}, they do not repair its
coordinatewise divisor provenance.
\end{proposition}

Thus the cubical construction improved the geometry of descent but left the
birth decision unchanged in substance.  The finite enumeration through
\(k=5\) confirmed exactly this product classification; it was a diagnostic
test of the factorization, not evidence for a genuine joint birth.

\section{Why a divisor shadow will still appear after forgetting provenance}

There is a separate point that must not be confused with the preceding
failure.  Suppose a future construction genuinely produces one birth
\(J_S\) for every finite set \(S\) of distinct single-prime births and has a
descent for every inclusion \(T\subseteq S\).  If one forgets all source
embeddings, attaching data, and filler spaces and remembers only the component
labels, the resulting descent poset is necessarily
\[
\{T:T\subseteq S\}.
\]
It is Boolean.  After posteriorly writing
\(S=\{p_1,\ldots,p_k\}\), the map
\[
T\longmapsto\prod_{p\in T}p
\]
identifies this Boolean shadow with the divisor lattice of the square-free
product.

Therefore no construction satisfying arbitrary-subcombination descent can
avoid having a divisor shadow after sufficient forgetting.  The correct
non-collapse demand is stronger and more precise:
\[
\boxed{
\text{the joint birth data must not be determined by that shadow.}
}
\]
The circle-cover proposal fails because the shadow is the whole object.  The
cubical proposal fails because its birth predicate is still determined by the
product shadow, while its extra topology is constant decoration.

\section{The one-candidate obstruction}

The preceding examples reveal a logical obstruction independent of the chosen
geometric presentation.

Let \(\mathcal X_k\) be a proposed space or category of \(k\)-birth
candidates, and let
\[
\rho_k:\mathcal X_k\longrightarrow
\operatorname{Sym}^k\bigl(\Pref(\mathscr R)_{\ge2}\bigr)
\]
record the unordered family of its one-dimensional path shadows.  Suppose
first that every fibre of \(\rho_k\) contains exactly one isomorphism class.

\begin{theorem}[One candidate per shadow forces unary factorization]
\label{thm:one-candidate}
Assume the desired exact classification is
\[
\Birth_k(X)=1
\quad\Longleftrightarrow\quad
\text{every path in }\rho_k(X)\text{ is a single-prime birth}.
\]
If \(\rho_k\) has one isomorphism class in each fibre, then \(\Birth_k\) is
completely determined by the unordered tuple of unary birth indicators.  In a
frame it is exactly
\[
\boxed{
\Birth_k(X)
=
\prod_{i=1}^k\Birth(H_i).
}
\]
Consequently no additional topology attached uniformly to the unique
candidate can be the deciding source of a genuinely mixed birth.
\end{theorem}

\begin{proof}
The displayed exact classification assigns value one precisely to the single
Boolean pattern \((1,\ldots,1)\) and value zero to every other pattern.
Therefore its value is the Boolean product of the unary indicators.  Since
there is only one candidate above each shadow, there is no remaining
within-fibre datum on which a different joint decision could depend.
\end{proof}

This theorem does not say that a non-product proof of the same numerical locus
is logically impossible.  It says that, with one candidate per shadow, such a
proof cannot produce additional joint birth data: the final verdict and the
candidate are already fixed by the unary shadow.  A constant cube, sphere,
relative class, or nerve placed over each tuple is therefore inert.

The Boolean support indicator is therefore the wrong place to look for
jointness.  If its support is exactly the set of all-prime tuples, that
indicator is unavoidably the product of the unary indicators.  A genuine
higher birth must retain a richer object, say a space or groupoid
\[
\mathfrak B_k(X),
\]
whose nonemptiness gives the Boolean support but whose topology and restriction
maps are not determined by that support.

\begin{corollary}[Necessary mixed-fibre test]
\label{cor:mixed-fibre}
For a topology-decisive joint theory, \(\rho_k\) must have nontrivial fibres,
and there must be candidates \(X\) and \(X'\) with the same one-dimensional
provenance shadows for which
\[
\mathfrak B_k(X)\not\simeq\mathfrak B_k(X')
\]
or for which their descent maps to lower birth objects are inequivalent.
Otherwise the enriched birth object also factors through the unary shadow,
possibly times an inert space depending only on \(k\).
\end{corollary}

Thus the source sought is necessarily a moduli problem, not a single canonical
product object.  For every finite prime multiset \(S\), the fibre over \(S\)
must contain at least one born object.  The Boolean statement that such a birth
exists may still factor through the unary indicators; the retained birth
\emph{object} must not.  Descent must act on the mixed fibre data as well as
forget a component.

\section{Coprime splitting in the multidirectional cover source}

The higher-rank continuation source from the source manuscript does provide
nontrivial moduli for several primitive directions of one fixed posterior
prime \(p\).  It is therefore necessary to test separately whether it couples
different primes.

Let \(G_d\) be its intrinsic abelian continuation group.  For a finite set
\(S\) of distinct posterior primes, a primitive direction over \(p\in S\) is
represented by a nonzero continuation character
\[
\chi_p:G_d\longrightarrow\mathbb F_p,
\]
up to multiplication by \(\mathbb F_p^\times\).  A simultaneous mixed-prime
quotient is represented by
\[
\chi_S:G_d\longrightarrow\prod_{p\in S}\mathbb F_p,
\qquad
\chi_S(g)=\bigl(\chi_p(g)\bigr)_{p\in S}.
\]

\begin{proposition}[Distinct-prime continuation moduli split]
\label{prop:coprime-splitting}
The moduli of mixed-prime continuation quotients in this abelian source is
\[
\boxed{
\mathfrak P_{d,S}
\cong
\prod_{p\in S}
\mathbb P\bigl((G_d/pG_d)^\vee\bigr).
}
\]
Every descent that removes one posterior prime is the projection deleting the
corresponding factor.  Hence this source contains no cross-prime joint moduli.
\end{proposition}

\begin{proof}
A homomorphism from \(G_d\) to the product is uniquely the tuple of its
coordinate characters.  If every character is nonzero, every coordinate
projection of the image is surjective.  The image order is then divisible by
every \(p\in S\), hence by their product; since it is a subgroup of a group of
that order, the simultaneous map is surjective.  Because the factors have
pairwise distinct prime orders, every automorphism of
\(\prod_{p\in S}C_p\) preserves each unique
\(p\)-primary subgroup.  Thus the change-of-generator action is exactly the
product of the scalar actions \(\mathbb F_p^\times\), with no automorphism
mixing two prime coordinates.  Quotienting the nonzero characters by these
independent scalar actions gives the displayed product.  Removing a factor
composes with the corresponding coordinate projection.
\end{proof}

Equivalently, a finite abelian group of square-free order is the canonical
product of its characteristic \(p\)-primary subgroups.  The projective
interaction topology available at one fixed characteristic does not cross
between distinct characteristics.  Therefore the source manuscript's
same-prime projective family cannot be promoted to an arbitrary distinct-prime
birth without adding a genuinely mixed source operation.

\section{Why nonabelian deck extensions do not solve the arbitrary-prime problem}

One might try to replace the abelian continuation quotient by a nonabelian
finite deck group whose simple shadows have different prime orders.  This can
produce mixed extensions for some pairs of primes, but not for arbitrary
pairs.

\begin{proposition}[Square-free two-prime obstruction]
\label{prop:pq-obstruction}
Let \(p<q\) be primes with \(p\nmid(q-1)\).  Every group of order \(pq\) is
cyclic.  In particular, every group of order \(15\) is \(C_{15}\).
\end{proposition}

\begin{proof}
Let \(G\) have order \(pq\).  The number \(n_q\) of Sylow \(q\)-subgroups
divides \(p\) and is congruent to one modulo \(q\), hence \(n_q=1\).  Thus the
\(q\)-subgroup is normal.  The number \(n_p\) divides \(q\) and is congruent
to one modulo \(p\), so it is either one or \(q\).  The second possibility
would imply \(q\equiv1\pmod p\), contrary to \(p\nmid(q-1)\).  Hence the
\(p\)-subgroup is also normal.  The two subgroups commute, have trivial
intersection, and generate \(G\), giving
\[
G\cong C_p\times C_q\cong C_{pq}.
\]
Taking \((p,q)=(3,5)\) gives the stated example.
\end{proof}

Therefore a finite regular-cover source whose joint object is supposed to
have exactly the prime deck factors in \(S\) cannot provide nontrivial mixed
extension topology for every \(S\).  Some finite prime sets force the deck
object to split canonically.  When nonabelian groups do exist for other prime
pairs, Proposition~\ref{prop:cover-subgroup} still makes intermediate-cover
descent a subgroup-series factorization.  Thus neither abelian nor nonabelian
finite deck extensions can be the uniform source of arbitrary finite-prime
joint births.

\section{The cross-primary coefficient obstruction}

Another tempting repair is to retain one homology class for each unary prime
birth and combine different classes by a product, smash product, or join.  For
distinct primes the mixed coefficient class disappears.

\begin{proposition}[No mixed class from primary tensor operations]
\label{prop:primary-tensor}
Let \(p\ne q\) be primes, let \(A\) be a \(p\)-primary abelian group, and let
\(B\) be a \(q\)-primary abelian group.  Then
\[
A\otimes_{\mathbb Z}B=0,
\qquad
\operatorname{Tor}_1^{\mathbb Z}(A,B)=0.
\]
Consequently, if the reduced homology of two unary birth carriers is purely
\(p\)-primary and \(q\)-primary respectively, the K\"unneth mixed terms in
their smash product vanish.  The same holds for the shifted smash model of
their join.
\end{proposition}

\begin{proof}
Every element of \(A\otimes B\) is killed by a power of \(p\) and by a power
of \(q\).  Coprimality gives a B\'ezout combination of these powers equal to
one, so the element is zero.  The standard cyclic calculation gives
\[
\operatorname{Tor}_1^{\mathbb Z}
(\mathbb Z/p^a,\mathbb Z/q^b)=0;
\]
passing through finite primary decompositions, and then filtered colimits when
needed, gives the assertion.  The topological statement follows from the
reduced K\"unneth theorem.
\end{proof}

Working over \(\mathbb F_p\) does not fix the problem: it privileges one
characteristic and erases torsion at every other prime.  Thus a mixed-prime
birth cannot be obtained merely by applying standard tensorial homology
operations to already separated unary prime classes.  Its coupling must be
carried by source geometry that exists integrally before primary
decomposition.

\section{Audit of the integral return-ribbon topology}

The return-ribbon sector of the source manuscript is a more serious candidate
than a product or coefficient construction.  Its critical incidence is the
literal cellular link \(L_s\) of a source hub, and its matching complex is
formed from actual overlaps of bounded-local rewrites.  It is integral and
source-native.  Nevertheless it does not retain the mixed-prime data required
here.

\begin{proposition}[The ribbon detector factors through total occurrence count]
\label{prop:ribbon-count}
At a retained stage \(H_s\), let \(N_s\) be the number of active nested return
occurrences.  The ribbon detector satisfies
\[
L_s\cong C_{N_s},
\qquad
\mathscr K_{\rm rib}(H_s)
\cong
\operatorname{Match}(C_{N_s})
\]
for \(N_s\ge3\), with the source manuscript's two-edge convention at
\(N_s=2\).  Posteriorly,
\[
N_s=\Omega(s)
\]
at every composite stage.  Hence the isomorphism and homotopy types of the
deciding complex factor through \(\Omega(s)\); they do not retain which prime
channels supplied the occurrences.
\end{proposition}

The provenance labels on radial germs are present in the full source state,
but the matching-complex birth detector uses only adjacency in the cyclic
link.  Relabelling the same cycle by different prime channels does not change
the detector.  For example,
\[
12=2^2\cdot3,
\qquad
30=2\cdot3\cdot5
\]
both have three occurrences and therefore the same unlabelled critical-cycle
matching type, despite having different prime multisets.

There is a second mismatch.  A face of
\(\operatorname{Match}(C_{N_s})\) is a family of pairwise disjoint local
critical rewrites.  Removing a matching edge does not remove a prime channel,
and maximal face chains do not terminate in the original single-prime birth
objects.  Thus the ribbon complex supplies a genuine topology for deciding the
unary prime/composite recursion, but it supplies neither mixed-prime fibres nor
the required prime-deletion descent.

The completed ribbon carrier confirms the same separation: posteriorly it is
the product \(\prod_p\mathbb Z_p\) of the primitive pro-return channels.
The local overlap detector enriches a composite stage, but it does not turn
that product of already retained channels into a new arbitrary-prime birth.

\section{Audit of the first genuinely coupled control}

The associahedral construction of Version~0.27 is source-native, but it cannot
serve as a finite-prime birth object in the present extraction.  For \(n\)
composable retained path pieces its critical shell is
\[
\operatorname{sd}(\partial K_n)\simeq S^{n-3}.
\]
This topology is determined by the number and order of the pieces, not by the
prime shadows carried by them.  Consequently every prime set of the same
cardinality receives the same associahedral shell.  More decisively, the cell
that closes this shell is explicitly the higher associativity-coherence cell.
It is therefore both arity-only relative to prime provenance and outside the
declared no-coherence scope.

Version~0.29 contains a more informative control.  Three interchange squares
share the same edge variables \(x,y,z\), so their joint filler is the coupled
space
\[
\operatorname{CubeFill}_{S}(a,b,c)
=
\{(x,y,z)\in S^3:[x,y]=a,\ [y,z]=b,\ [z,x]=c\},
\]
not the product of three independent square-filler sets.  In the displayed
\(A_5\) example every one-face and every two-face subsystem is fillable, while
the full filler is empty.  Its existence nerve is therefore
\(\partial\Delta^2\simeq S^1\).  A contrasting \(A_5\) packet has the same
one-face fillability property and a full filler, hence a contractible nerve.
Thus shared source variables can make global topology affect a birth decision
without reducing to coordinatewise Boolean data.

This is the first structural analogue in the source manuscript of the mixed-
fibre behaviour required here, but it is not an instance of prime-combination
birth.  It fails the requested specification in four separate ways:

\begin{enumerate}[label=(\roman*)]
\item the coefficient atom \(A_5\) is selected for the control rather than born
from the one-direction prime source;
\item the construction compares nonlinear commutator defects, not arbitrary
finite sets of distinct prime births;
\item its deciding object is explicitly a coherent globalization or homotopy-
limit filler;
\item its face restrictions do not delete prime channels and do not terminate
in the original unary regular-atlas births.
\end{enumerate}

It would therefore be incorrect to rename the shared-edge \(A_5\) packet as a
\(k\)-prime birth.  What the example proves is narrower but useful: any genuine
replacement must couple raw source variables before unary classification, and
two candidates with the same unary shadows must still be able to have different
joint topological fibres.  The source manuscript realizes this phenomenon only
in the nonlinear coherence sector that has been excluded from the present
paper.

\section{A finite-packet trial: one path with overlapping returns}

The preceding audit leaves open a construction that is still genuinely
one-path.  Let \(L\) be a finite directed atomic interval retained by the
original homogeneous source, and let
\[
q:L\longrightarrow Y
\]
be a cellular path trace which may identify vertices but does not collapse the
interior of an atomic edge.  A retained \emph{return support} is a nonempty
subinterval \(J=[a,b]\subseteq L\) with \(q(a)=q(b)\).  Its unary normalization
is the cut-open ordered path
\[
N(J):=J.
\]
Thus several returns may share literal atomic edges of one chronological path;
they are not separate born paths subsequently glued together.  The quotient
trace, the supports, and their edge intersections are retained before any
posterior rank or prime name is assigned.

For a finite return packet
\(\mathcal J=\{J_1,\ldots,J_k\}\), let
\(\Gamma(\mathcal J)\) be the graph whose vertices are the retained returns and
in which
\[
J_i\sim J_j
\quad\Longleftrightarrow\quad
J_i\cap J_j\text{ contains an atomic edge of }L.
\]
Call the packet \emph{source-separable} if it has a nontrivial partition
\(\mathcal J=\mathcal A\sqcup\mathcal C\) such that no support in
\(\mathcal A\) shares an atomic edge with a support in \(\mathcal C\).
Otherwise call it \emph{source-indecomposable}.

This is not an extra connectivity convention: it is exactly the obstruction to
writing the retained return data as two edge-disjoint path batches.  Directly
from the definitions,
\[
\mathcal J\text{ is source-indecomposable}
\quad\Longleftrightarrow\quad
\Gamma(\mathcal J)\text{ is connected}
\quad\Longleftrightarrow\quad
\widetilde H_0(\Gamma(\mathcal J);\mathbb Z)=0.
\]
The last condition is ordinary topology of the actual source-support
intersection graph; it is not a coherence or filler condition.

The finite-packet provenance record is
\[
\Prov_{\cap}(q;\mathcal J)
=
\left(
q:L\to Y,
\mathcal J,
\bigl(J_i\cap J_j\bigr)_{i,j},
\bigl(\RegAtlas(N(J_i);\Pi)\bigr)_{i,\Pi}
\right),
\]
where \(\Pi\) ranges through the earlier unary birth ledger appropriate to the
normalized support.  Define the trial packet predicate by
\[
\boxed{
\Birth_{\cap}(q;\mathcal J)=1
\quad\Longleftrightarrow\quad
\begin{cases}
\Gamma(\mathcal J)\text{ is connected},\\
\RegAtlas(N(J_i);\Pi)=\varnothing
\text{ for every }i\text{ and every earlier unary birth }\Pi.
\end{cases}}
\]
The candidate packet and all of its intersections are source data.  The unary
ledger is used only to apply the already established regular-atlas test to its
cut-open faces; no completed unary birth object is glued into the packet.

Applying posterior edge count only after this predicate has been formed gives
the exact finite-packet classification
\[
\Birth_{\cap}(q;\mathcal J)=1
\quad\Longleftrightarrow\quad
\Gamma(\mathcal J)\text{ is connected and every }
\rk(N(J_i))\text{ is prime}.
\]
Indeed, source-indecomposability is the first condition by definition, while
the single-prime theorem applied separately to each normalized return gives the
second.  The Boolean support is necessarily the all-prime locus; the additional
joint datum is the retained intersection topology over that locus.

This additional datum passes the mixed-fibre test.  On one ambient atomic path
take return supports of posterior lengths \(2\) and \(3\).  The two placements
\[
[0,2],\ [1,4]
\qquad\text{and}\qquad
[0,2],\ [3,6]
\]
have the same two unary shadows.  In the first placement the supports share an
atomic edge and the packet is born; in the second they are source-separated and
there is no joint birth.  Hence the packet verdict is not recoverable from the
two unary provenance shadows.

There is also a genuine descent tower.  A deletion
\[
d_i(q;\mathcal J):=(q;\mathcal J\setminus\{J_i\})
\]
is admissible when the remaining intersection graph is connected.  Every
finite connected graph with at least two vertices has such a deletion: choose
a spanning tree and remove one of its leaves.  Repeating this operation gives
\[
(q;\mathcal J_k)
\longrightarrow
(q;\mathcal J_{k-1})
\longrightarrow\cdots\longrightarrow
(q;\mathcal J_1),
\qquad |\mathcal J_r|=r,
\]
with every intermediate packet still source-indecomposable.  If the top packet
is born, every normalized return in it is unary-prime, so every stage of the
tower is born and the final object normalizes to the original single-prime
path \(H_p\).

This descent is not the divisor lattice.  Its admissible faces are the
connected induced subpackets of \(\Gamma(\mathcal J)\), and that poset changes
when the overlap pattern changes while all unary ranks remain fixed.  For
example, three returns of the same unary ranks can have intersection graph
\(P_3\) or \(K_3\).  The subpacket consisting of the two end vertices is absent
in the first connected-face poset and present in the second.  No divisor data
of the unary ranks can distinguish these cases.

At finite-packet level this construction therefore satisfies the three tests:
one connected path trace precedes its unary shadows, the descent terminates in
the original unary mechanism, and shared-support topology is not determined by
a divisor or product shadow.  At this point the only missing obligation is a
single unlabelled source which exposes the return packets without receiving a
chosen list of endpoint identifications.  The next section supplies that
source.

\section{The universal unlabelled path-history source}

The missing source can be supplied without choosing a prime list or a return
schedule.  For each finite directed atomic interval \(L\), consider all vertex
equivalence relations \(\sim\) satisfying
\[
v\not\sim w
\qquad
\text{whenever \(v,w\) are joined by one atomic edge.}
\]
The quotient map \(q_\sim:L\to L/{\sim}\) is one finite path history: the path
may revisit an earlier state, but no atomic move is collapsed.  Equivalence
classes have no labels.  Two histories are identified when an isomorphism of
directed atomic intervals carries one equivalence relation to the other.

Let \(\mathscr R_{\rm ret}\) be the rooted system of all these finite histories.
An atomic extension appends one edge and one terminal vertex.  The new vertex
is either a new equivalence class or is placed in one already exposed class
other than the class of the current endpoint; old classes are never merged.
This is one uniform path-extension rule.  At each finite stage it has only
finitely many possible extensions, and none is selected using a number-theoretic
predicate.

\begin{proposition}[Canonical normal form of the return source]
\label{prop:return-normal-form}
Length-\(n\) objects of \(\mathscr R_{\rm ret}\) are in bijection with words
\[
c_0c_1\cdots c_n
\]
satisfying
\[
c_i\in\mathbb N,
\qquad
c_0=0,
\qquad
c_i\le 1+\max_{j<i}c_j,
\qquad
c_i\ne c_{i-1}.
\]
The symbols are normalized by order of first occurrence.  Under this
bijection, adjoining the next unused symbol is the all-new atomic extension,
while adjoining any allowed old symbol is a return.  Hence the stated source
rule produces every unlabelled non-edge-collapsing path history exactly once
up to isomorphism.
\end{proposition}

\begin{proof}
Traverse a finite history from its root and number its equivalence classes by
the order in which they first appear.  A new class must receive the next unused
number, whereas a revisit receives its existing number.  Adjacent vertices
cannot receive the same number because no atomic edge is collapsed.  This
gives the displayed conditions and is unchanged by unlabelled history
isomorphism.

Conversely, such a word defines an equivalence relation by
\(i\sim j\) exactly when \(c_i=c_j\).  The last inequality prevents an atomic
edge from collapsing, and normalization by first occurrence makes the word
unique.  The two constructions are inverse, and deleting the terminal symbol
is exactly finite-history restriction.
\end{proof}

This proposition isolates the one genuinely new finite-source permission:
an atomic path may revisit a noncurrent exposed state.  The unary mother does
not force that permission.  Once it is granted, however, no further selection
of return patterns, prime types, spacings, or packet arities is made;
\(\mathscr R_{\rm ret}\) is the exhaustive unlabelled return extension.

For each atomic depth \(n\), let \(\mathscr H_n\) be the finite discrete space
of length-\(n\) path histories, and let
\[
r_{n+1,n}:\mathscr H_{n+1}\longrightarrow\mathscr H_n
\]
forget the last atomic edge and terminal occurrence.  Every history has an
all-new extension, so these restriction maps are surjective.  Define the
infinite-history completion by
\[
\boxed{
\widehat{\mathscr R}_{\rm ret}
:=
\varprojlim_n(\mathscr H_n,r_{n+1,n}).}
\]
Its basic cylinder \(C(h)\), for \(h\in\mathscr H_n\), consists of the infinite
histories whose length-\(n\) restriction is \(h\).  The cylinders are clopen
and form a basis.  Hence \(\widehat{\mathscr R}_{\rm ret}\) is a nonempty
compact zero-dimensional space.  This is ordinary inverse-limit topology on
path histories; no measure or coherence data are added.

The branch on which every new vertex forms a new class is canonically the
original atomic successor source
\[
o\longrightarrow x_1\longrightarrow x_2\longrightarrow\cdots.
\]
It will be called the \emph{unary spine}.  Its regular-atlas ledger and its
single-prime births are therefore unchanged.

Whenever a path extension revisits an exposed class, let \(a\) be the most
recent earlier occurrence of that class and let \(b\) be the new occurrence.
The interval \([a,b]\subset L\) is then an \emph{elementary return support}.
It is produced by the extension history itself; there is no separately supplied
list of endpoint identifications.

For the finite-combination theorem it suffices to use the \emph{simple-return
sector}: every recurrent equivalence class occurs exactly twice, and the
normalized elementary return supports are pairwise nonisomorphic as ordered
atomic paths.  Both conditions are properties of the retained unlabelled path
history.  Posteriorly, the second says that the return ranks are distinct.  The
complete return packet \(\mathcal J(X)\) of a simple-return history \(X\) is the
set of all its two-occurrence classes; it is not a chosen subfamily.

The packet predicate of the preceding section now becomes a source predicate:
\[
\boxed{
\Birth_{\rm ret}(X)=1
\quad\Longleftrightarrow\quad
\Gamma(\mathcal J(X))\text{ is connected and every normalized return has}
\text{ empty regular-atlas ledger}.}
\]
All terms on the right are retained before the posterior ranks are named.

\begin{theorem}[Arbitrary finite distinct-prime return birth]
\label{thm:return-combination}
The posterior shadow of \(\Birth_{\rm ret}\) is exactly the class of nonempty
finite sets of distinct primes.  More precisely:
\begin{enumerate}[label=(\roman*)]
\item every born simple-return packet has pairwise distinct prime unary ranks;
\item for every finite set \(S=\{p_1,\ldots,p_k\}\) of distinct primes there is
a born history \(X_S\in\mathscr R_{\rm ret}\) whose normalized return ranks are
exactly \(S\);
\item every such born history admits a descent tower through born connected
return packets to one original unary birth.
\end{enumerate}
\end{theorem}

\begin{proof}
For (i), empty regular-atlas ledgers imply by the single-prime theorem that all
normalized return ranks are prime.  Pairwise nonisomorphism in the simple-return
sector makes those ranks distinct.

For (ii), write the elements of \(S\) in increasing posterior order
\[
p_1<p_2<\cdots<p_k.
\]
This ordering is used only to verify existence in the numerical shadow.  Set
\[
a_1=0,
\qquad b_1=p_1,
\qquad
a_i=b_{i-1}-1,
\qquad
b_i=a_i+p_i
\quad(2\le i\le k).
\]
On the atomic interval \(L_{b_k}\), make \(\{a_i,b_i\}\) a two-element
equivalence class for every \(i\), and leave every other vertex in a singleton
class.

All \(2k\) displayed endpoints are distinct.  Indeed,
\(0<a_2<b_1\), while for \(i\ge3\) the distinctness of the primes gives
\(p_{i-1}\ge3\), hence
\[
b_{i-2}<a_i=b_{i-1}-1<b_{i-1}.
\]
Also \(b_i>b_{i-1}\), so no later endpoint collides with an earlier one.  Since
every \(p_i\ge2\), no paired endpoints are adjacent.  The resulting equivalence
relation is therefore an object of \(\mathscr R_{\rm ret}\), and its elementary
returns are exactly
\[
J_i=[a_i,b_i],
\qquad \rk(N(J_i))=p_i.
\]

Consecutive supports meet in precisely the atomic edge
\([b_{i-1}-1,b_{i-1}]\).  If \(j\le i-2\), then
\(b_j\le b_{i-2}<a_i\), so \(J_j\) and \(J_i\) are disjoint.  Consequently
\[
\Gamma(\mathcal J(X_S))\cong P_k.
\]
The packet is source-indecomposable, and every one of its normalized returns is
a unary prime birth.  Thus \(X_S\) is born.

For (iii), split the endpoint pair of \(J_k\), then that of \(J_{k-1}\), and so
on.  After each split the remaining intersection graph is
\(P_{k-1},P_{k-2},\ldots,P_1\), hence is connected.  All remaining atlas
ledgers stay empty.  The last support normalizes to \(H_{p_1}\), an original
unary birth.  For a general born packet, the spanning-tree leaf argument from
the preceding section supplies the same kind of connected-deletion tower.
\end{proof}

The descent arrow is internal to the normalized source presentation.  Splitting
a two-occurrence equivalence class into two singleton classes removes one
return while leaving the atomic interval and every other equivalence class
unchanged.  On quotient graphs the corresponding continuous map points in the
opposite direction, from the split history to the pinched history; the descent
category records the refinement of retained endpoint identifications.  No
intermediate cover or integer factor is involved.

The source is branching only in the space of possible path histories.  Each
candidate \(X\), every return packet, and every descent tower still belongs to
one chronological path.  Thus \(\mathscr R_{\rm ret}\) is not a product of
independent paths or a batch of completed unary births.  Branching is necessary
only because incompatible choices ``visit a new state'' and ``return to this
earlier state'' cannot both be the next step of one deterministic history.


\section{The source-transitive all-prime boundary birth}

The universal source also has a boundary of infinite path histories.  An object
of \(\widehat{\mathscr R}_{\rm ret}\) is one directed atomic ray \(L_\infty\)
with a vertex equivalence relation whose restriction to every finite interval
is an object of \(\mathscr R_{\rm ret}\).  This completion is defined before
the birth ledger and adds no local extension rule.

The primitive finite types needed at the boundary can be recognized without
naming primes.  Call a finite nonunit retained path \(Q\)
\emph{atlas-irreducible} when
\[
\RegAtlas(Q;R)=\varnothing
\qquad
\text{for every nonunit path type }R\subsetneq Q.
\]
Let \(\operatorname{Irr}_{\rm at}(\mathscr R)\) be the family of these raw path
types.

\begin{lemma}[Atlas-irreducibles are the unary births]
\label{lem:atlas-irreducible}
A finite nonunit path is atlas-irreducible exactly when it is born by the
original unary mother rule.  Consequently the posterior ranks of
\(\operatorname{Irr}_{\rm at}(\mathscr R)\) are exactly the primes.
\end{lemma}

\begin{proof}
If a path is inherited by the unary mother, its regular atlas by an earlier
primitive path proves that it is not atlas-irreducible.  Conversely, suppose
\(Q\) has a regular atlas by a shorter nonunit path \(R\).  Posterior edge count
gives
\[
\rk(Q)=m\,\rk(R),
\qquad m\ge2,
\]
so \(\rk(Q)\) is composite.  By the single-prime theorem \(Q\) is not a unary
birth.  The final assertion follows from the same theorem.  Arithmetic is used
only in this external identification of the already defined path property.
\end{proof}

There is a canonical operation on the completed source: forget the first
atomic step and re-root the remaining ray.  In terms of the vertex equivalence
relation of \(X\), define the \emph{path shift} \(T\) by
\[
i\sim_{TX}j
\quad\Longleftrightarrow\quad
i+1\sim_Xj+1.
\]
It is continuous.  Indeed, the inverse image of a finite-prefix cylinder is a
finite union of finite-prefix cylinders.

The closed source generated by a boundary history is therefore canonically
\[
\operatorname{Hull}_T(X)
:=
\overline{\{T^mX:m\ge0\}}.
\]
It is the smallest closed forward-shift-invariant subsource containing \(X\).

For \(m,n\ge0\), write
\[
W_{m,n}(X):=(T^mX)|_{L_n}
\]
for the length-\(n\) window beginning at time \(m\), and define the finite
language of \(X\) by
\[
\mathcal L(X)
:=
\{[W_{m,n}(X)]:m,n\ge0\}.
\]
This is a language of unlabelled path histories, not a list of ranks.

The finite mother alone cannot logically decide what happens on the new
inverse-limit boundary.

\begin{proposition}[Finite provenance does not determine a boundary verdict]
\label{prop:finite-not-boundary}
The birth predicates on the finite spaces \(\mathscr H_n\), even together
with path-history isomorphism invariance, do not determine a unique birth
predicate on \(\widehat{\mathscr R}_{\rm ret}\).
\end{proposition}

\begin{proof}
Keep every finite verdict unchanged.  One extension may declare every
infinite history born.  Another may declare only histories with full orbit
hull born.  Both agree with all finite verdicts and are invariant under
path-history isomorphism, but they disagree, for example, on the all-new unary
spine.
Thus some boundary inheritance principle is unavoidable; it cannot be proved
from finite verdicts alone.
\end{proof}

We use the following minimal topological continuation of the word
``inherited.''

\medskip
\noindent\emph{Generated-subsource inheritance principle.}
An infinite history \(X\) is inherited at the boundary exactly when the
smallest closed forward-shift-invariant subsource containing it is proper.
It is born exactly when it is not so inherited.
\medskip

There is no choice of subsource hidden in this principle.  The smallest one
is uniquely \(\operatorname{Hull}_T(X)\): every closed forward-invariant set
containing \(X\) contains the forward orbit and hence its closure, while the
orbit closure itself has those properties.  Therefore the following boundary
rule is forced once the generated-subsource inheritance principle is adopted.

\begin{definition}[Source-transitive boundary birth]
\label{def:source-transitive-birth}
An infinite path history \(X\in\widehat{\mathscr R}_{\rm ret}\) is
\emph{born at the source boundary} when
\[
\boxed{
\Birth_\infty(X)=1
\quad\Longleftrightarrow\quad
\operatorname{Hull}_T(X)
=
\widehat{\mathscr R}_{\rm ret}.}
\]
Equivalently, \(X\) is source-transitive.
\end{definition}

If the orbit closure is proper, then \(X\) is inherited from that canonical
proper subsource.  If the orbit is dense, no proper closed shift-invariant
subsource contains \(X\).  Thus the distinction is made by the topology and
the source shift alone; no prime predicate, spacing rule, probability law, or
coherence datum enters the definition.

\begin{lemma}[Finite-language criterion]
\label{lem:finite-language-criterion}
For every \(X\in\widehat{\mathscr R}_{\rm ret}\),
\[
\boxed{
\operatorname{Hull}_T(X)
=
\{Y\in\widehat{\mathscr R}_{\rm ret}:
  \mathcal L(Y)\subseteq\mathcal L(X)\}.}
\]
Consequently the following are equivalent:
\begin{enumerate}[label=(\roman*)]
\item \(X\) is a source-transitive boundary birth;
\item every finite history type occurs as a window of \(X\);
\item
\[
\mathcal L(X)=\bigcup_{n\ge0}\mathscr H_n.
\]
\end{enumerate}
\end{lemma}

\begin{proof}
Suppose \(Y\in\operatorname{Hull}_T(X)\).  Every finite prefix of every
\(T^rY\) is then a limit of prefixes of shifted copies of \(X\).  Since the
finite history spaces are discrete, that prefix actually occurs as a window
of \(X\).  Thus \(\mathcal L(Y)\subseteq\mathcal L(X)\).

Conversely, suppose \(\mathcal L(Y)\subseteq\mathcal L(X)\).  For each \(n\),
the prefix \(Y|_{L_n}\) occurs as some window \(W_{m_n,n}(X)\).  Hence
\(T^{m_n}X\) lies in the length-\(n\) cylinder of \(Y\).  These cylinders form
a neighborhood basis at \(Y\), so \(Y\in\operatorname{Hull}_T(X)\).  This
proves the displayed identity.

Finally, the orbit of \(X\) meets the cylinder specified by
\(F\in\mathscr H_n\) exactly when \(W_{m,n}(X)\cong F\) for some \(m\).
Since the finite-prefix cylinders form a basis, meeting all of them is
equivalent both to density and to
\(\mathcal L(X)=\bigcup_n\mathscr H_n\).
\end{proof}

\begin{theorem}[Existence of source-transitive path histories]
\label{thm:transitive-existence}
Source-transitive boundary births exist in
\(\widehat{\mathscr R}_{\rm ret}\).
\end{theorem}

\begin{proof}
The family \(\bigcup_{n\ge0}\mathscr H_n\) is countable because every
\(\mathscr H_n\) is finite.  Choose representatives
\(F_1,F_2,\ldots\) in which every finite history type occurs.  Place these
finite words consecutively on one atomic ray.  Give the equivalence classes
used by different words disjoint names and insert a fresh singleton vertex
between consecutive words.  Inside each word retain exactly its given
equivalence relation; make no identification across two words.

Adjacent global vertices remain inequivalent, so the resulting infinite
history \(X\) belongs to \(\widehat{\mathscr R}_{\rm ret}\).  Each \(F_i\)
occurs as the window beginning at its first vertex.  Hence
\(\mathcal L(X)=\bigcup_n\mathscr H_n\), and
Lemma~\ref{lem:finite-language-criterion} makes \(X\) a boundary birth.
The enumeration is only an existence proof.  It is not part of the birth
predicate, and different dense histories need not be isomorphic.
\end{proof}

For a finite window \(F\), let \(\mathcal J(F)\) be its elementary return
supports.  Define the primitive normalized-return spectrum visible in an
infinite history by
\[
\operatorname{Spec}_{\rm prim}(X)
:=
\left\{
[N(J)]:
\begin{array}{l}
F\in\mathcal L(X),\ J\in\mathcal J(F),\\
N(J)\in\operatorname{Irr}_{\rm at}(\mathscr R)
\end{array}
\right\}.
\]
The filter here is the raw regular-atlas test of
Lemma~\ref{lem:atlas-irreducible}.  Composite returns may occur in the
universal history, just as composite stages already occur in the original
unary source; they are inherited data, not primitive births.

\begin{theorem}[One-path realization of the full prime distribution]
\label{thm:all-prime-return}
Let \(X\) be any source-transitive boundary birth.  Then:
\begin{enumerate}[label=(\roman*)]
\item
\[
\operatorname{Spec}_{\rm prim}(X)
=
\operatorname{Irr}_{\rm at}(\mathscr R);
\]
hence its posterior primitive rank spectrum is exactly \(\mathbb P\);
\item every born finite return packet occurs as a window in this same path
history \(X\);
\item for every nonempty finite set \(S\subset\mathbb P\), some window of
\(X\) is a joint birth with posterior spectrum \(S\), and that window has a
connected-deletion descent tower ending at each chosen unary member.
\end{enumerate}
\end{theorem}

\begin{proof}
For an atlas-irreducible finite path \(Q\), form the finite history having
fresh internal vertices and one endpoint class occurring at the two ends of
\(Q\).  It has one elementary return whose normalization is \(Q\).  This
finite history occurs as a window of \(X\) by
Lemma~\ref{lem:finite-language-criterion}.  Therefore every member of
\(\operatorname{Irr}_{\rm at}(\mathscr R)\) lies in
\(\operatorname{Spec}_{\rm prim}(X)\).  The reverse inclusion is part of the
definition of the primitive spectrum.  Lemma~\ref{lem:atlas-irreducible}
then gives the posterior equality with \(\mathbb P\).

Every finite born packet is itself an element of some \(\mathscr H_n\), so it
also occurs as a window of \(X\).  Apply this to the finite packet \(X_S\)
constructed in Theorem~\ref{thm:return-combination}.  Its return-support
graph is connected and its primitive return ranks are exactly \(S\).
The connected-deletion theorem, with a spanning tree rooted at any selected
return, supplies a descent tower inside that window to the selected unary
birth.
\end{proof}

\begin{proposition}[Boundary-to-unary descent tower]
\label{prop:boundary-unary-descent}
Let \(X\) be a source-transitive boundary birth, let
\(\varnothing\ne S\subset\mathbb P\) be finite, and choose \(p\in S\).
There are \(m,n\) and a tower
\[
X
\xrightarrow{\ T^m\ }
T^mX
\xrightarrow{\ \rho_n\ }
F_S
\xrightarrow{\ d_{J_1}\ }
F_{S_1}
\longrightarrow\cdots\longrightarrow
F_{\{p\}}
\xrightarrow{\ N\ }
H_p,
\]
where \(\rho_n\) is finite-prefix restriction, every infinite term is a
boundary birth, every finite packet term is a finite return birth, every
\(d_{J_i}\) splits one endpoint class, and \(H_p\) is the original unary
birth of posterior rank \(p\).
\end{proposition}

\begin{proof}
Every nonempty cylinder in \(\widehat{\mathscr R}_{\rm ret}\) has infinitely
many extensions: after its prescribed prefix one may postpone a new return by
arbitrarily many fresh vertices.  Hence the completion has no isolated point.
Removing finitely many points from a dense orbit therefore leaves a dense
set, so every tail \(\{T^{m+r}X:r\ge0\}\) is dense.  Thus \(T^mX\) is again a
boundary birth.

By Theorem~\ref{thm:all-prime-return}, choose a window
\(W_{m,n}(X)\cong F_S\) which is a born connected packet with spectrum \(S\).
The map \(\rho_n\) selects precisely this prefix of \(T^mX\).  Root a spanning
tree of \(\Gamma(\mathcal J(F_S))\) at the return of rank \(p\), and repeatedly
split an endpoint class belonging to a nonroot leaf.  The remaining induced
support graph stays connected at every step, while the atlas ledger of every
remaining return is unchanged.  Hence every \(F_{S_i}\) is born.  The last
packet has one return, and its normalization is the original unary path
\(H_p\).  All arrows are literal operations on one path history; no comparison
or coherence map is inserted.
\end{proof}

\begin{corollary}[The boundary birth does not collapse mixed fibres]
\label{cor:boundary-mixed-fibres}
Let \(X\) be source-transitive and let
\(S\subset\mathbb P\) be finite with \(|S|\ge2\).  Then \(X\) has two finite
windows \(F_S^+\) and \(F_S^-\) with the same primitive posterior spectrum
\(S\), while
\[
\Birth_{\rm ret}(F_S^+)=1,
\qquad
\Birth_{\rm ret}(F_S^-)=0.
\]
Consequently the finite birth verdict inside \(X\) does not factor through
the prime set, its product, or its divisor poset.
\end{corollary}

\begin{proof}
Take \(F_S^+\) to be the connected packet from
Theorem~\ref{thm:return-combination}.  For \(F_S^-\), place normalized returns
of the same primitive types consecutively on a finite atomic interval, using
fresh endpoint classes and at least one atomic edge between successive
supports.  Its return-support graph is discrete, so it is not
source-indecomposable, although its normalized primitive returns have the same
posterior ranks \(S\).  Both finite histories occur as windows of \(X\) by
Lemma~\ref{lem:finite-language-criterion}.
\end{proof}

Thus a single \(X\) contains all finite prime combinations, not merely all
unary prime types.  It does so through different finite windows of one
chronological path.  No fixed overlap graph is imposed at infinity: the
actual finite overlap graph belongs to the window in which a packet occurs.
In particular, equal unary shadows with different mixed fibres both occur,
so the global construction retains rather than collapses the finite joint
geometry.

The pointed realization provenance keeps the carrier and the positions of all
of its windows:
\[
\Prov^\bullet_\infty(X)
=
\left(
 X,\,
 \operatorname{Hull}_T(X),\,
 \mathcal L(X),\,
 \bigl(
   W_{m,n}(X),\,
   \mathcal J(W_{m,n}(X)),\,
   \Gamma(\mathcal J(W_{m,n}(X)))
 \bigr)_{m,n\ge0}
\right).
\]
This pointed record distinguishes nonisomorphic dense histories.

The birth mother does not need those time coordinates.  Its source-level
provenance retains the generated subsource, the finite language, and all local
packet types, but forgets where and how often a type occurs:
\[
\Prov^{\rm src}_\infty(X)
=
\left(
 \operatorname{Hull}_T(X),\,
 \mathcal L(X),\,
 \left\{
 \left(
   F,\,
   \mathcal J(F),\,
   \Gamma(\mathcal J(F)),\,
   \bigl(\RegAtlas(N(J);\Pi)\bigr)_{J,\Pi}
 \right):
 [F]\in\mathcal L(X)
 \right\}
\right).
\]
The packet and atlas data in the last component are determined by the finite
history type \(F\).  Thus equality of source-level provenance is equivalent to
equality of finite languages.  Define
\[
X\equiv_{\Prov}Y
\quad\Longleftrightarrow\quad
\mathcal L(X)=\mathcal L(Y).
\]
By Lemma~\ref{lem:finite-language-criterion}, this is also equivalent to
\(\operatorname{Hull}_T(X)=\operatorname{Hull}_T(Y)\).

\begin{theorem}[Uniqueness of the all-prime provenance class]
\label{thm:unique-provenance-class}
Let
\[
\mathfrak B_\infty
:=
\{X\in\widehat{\mathscr R}_{\rm ret}:\Birth_\infty(X)=1\}.
\]
Then
\[
\boxed{
\mathfrak B_\infty/{\equiv_{\Prov}}
\text{ consists of exactly one class}.}
\]
That class has generated subsource
\(\widehat{\mathscr R}_{\rm ret}\), finite language
\(\bigcup_n\mathscr H_n\), and posterior primitive spectrum \(\mathbb P\).
\end{theorem}

\begin{proof}
If \(X,Y\in\mathfrak B_\infty\), then the finite-language criterion gives
\[
\mathcal L(X)
=
\bigcup_n\mathscr H_n
=
\mathcal L(Y),
\]
so \(X\equiv_{\Prov}Y\).  Existence follows from
Theorem~\ref{thm:transitive-existence}.  Conversely, any history
provenance-equivalent to one member of \(\mathfrak B_\infty\) has the full
finite language and therefore is itself a boundary birth.  The statements
about the generated source and primitive spectrum follow from
Definition~\ref{def:source-transitive-birth} and
Theorem~\ref{thm:all-prime-return}.
\end{proof}

This quotient does not collapse the mixed geometry used by the theory: every
finite history type \(F\), its actual return-support graph, and both verdicts
in Corollary~\ref{cor:boundary-mixed-fibres} remain in the common language.
It forgets only the root-relative occurrence positions and multiplicities.
The equality \(\mathcal L(X)=\bigcup_n\mathscr H_n\) is not certified by one
finite prefix; it is a genuine boundary property.  Each individual
requirement, however, is witnessed by a finite window, and every finite
descent is a literal restriction and endpoint-class splitting inside that
window.  No comparison map between alternative descents is asserted or
needed.

\section{The global descent tradeoff}

The source-transitive object realizes all finite combinations, but it does not
do so by choosing one global representative of each prime type and declaring
all subsets of those representatives to be faces.  That stronger one-copy
model would recover the divisor shadow.

\begin{proposition}[Universal one-copy descent forces the Boolean shadow]
\label{prop:global-boolean-tradeoff}
Let \(Y\) be a return object having exactly one return \(J_p\) for every unary
prime type, and declare a finite subpacket jointly admissible exactly when its
induced return-support intersection graph is connected.  If every two-prime
subpacket \(\{J_p,J_q\}\) is an admissible descent face of \(Y\), then
\[
\Gamma(\mathcal J(Y))\cong K_{\mathbb P}.
\]
Consequently every nonempty finite prime subset is an admissible face and the
finite-face poset is the Boolean poset
\[
\{S\subset\mathbb P:0<|S|<\infty\}.
\]
After the posterior map \(S\mapsto\prod_{p\in S}p\), this is the poset of
square-free divisors ordered by divisibility.
\end{proposition}

\begin{proof}
A graph on two vertices is connected exactly when its two vertices are joined
by an edge.  Hence admissibility of every pair \(\{J_p,J_q\}\) forces
\(J_p\sim J_q\) for every distinct \(p,q\), so the full intersection graph is
complete.  Every nonempty finite induced subgraph of a complete graph is
connected, giving all finite subsets and all subset inclusions.  Unique prime
factorization identifies that Boolean poset with the stated square-free
divisor poset only after posterior ranks have been read.
\end{proof}

The proposition diagnoses the discarded exact-one-copy model; it does not
apply to Theorem~\ref{thm:all-prime-return}.  A source-transitive history
contains many occurrences of a primitive path type.  A finite set \(S\) is
realized by one connected window, and another set may be realized by a
different window using different occurrences.  There is therefore no
prime-indexed global vertex set whose entire subset lattice is being declared
to be the descent geometry.  Each finite descent tower is the connected-face
poset of its actual window and depends on the return-overlap incidence of that
window, not on divisibility.

\section{Boundary of the construction}

The finite theorem resolves arbitrary finite sets of distinct primes inside
one uniform unlabelled path-history source.  Theorem~\ref{thm:all-prime-return}
then strengthens this in the requested second sense: a single infinite path
history contains every such finite joint birth as an actual finite window,
as well as every unary prime birth.

The boundary object is not required to be simple-return.  Repetition is
essential here: the same primitive type can occur in many windows, allowing
incompatible mixed fibres to coexist at different locations of one
chronological path.  Composite returns and equivalence classes with further
visits may also occur, but the raw atlas test separates inherited returns from
the primitive spectrum.  The distinct-prime packet theorem itself remains in
the simple-return sector, where its classification is exact.

There is no canonical spacing, overlap graph, or enumeration in the boundary
birth rule.  The shift-orbit condition selects a generally non-singleton
moduli of source-transitive histories.  The enumeration in
Theorem~\ref{thm:transitive-existence} proves that this locus is nonempty but
does not become source-provenance data.  Theorem~\ref{thm:unique-provenance-class}
shows that the entire moduli is one provenance class: what is invariant is
the full finite language, in which every permitted finite source history
occurs.

The earlier negative audits retain their force.  Intermediate covers give a
divisor poset, cubical atlases give a product of divisor posets, abelian
mixed-prime covers split by primary type, ribbon matching topology forgets the
prime labels, and the \(A_5\) shared-edge control belongs to the excluded
coherence sector.  The return-history source avoids those failures because
its mixed datum is the actual atomic-edge intersection pattern of elementary
returns before unary classification.  The boundary step adds only the
canonical shift and the topology already carried by the inverse-limit source.

\section{Closure of the proof}

The construction now closes at four distinct levels.

First, after adjoining the single permission to revisit an earlier state, the
source closes: Proposition~\ref{prop:return-normal-form} proves that the
new-class/old-class alternatives generate every unlabelled
non-edge-collapsing path history exactly once, and the inverse-limit completion
contains exactly the compatible infinite histories.

Second, finite joint birth closes: connected return support is a source
property, the empty unary atlas ledger is the primitive-face property, every
finite distinct-prime set has such a packet, and every born packet has a
connected-deletion tower to any selected unary member.

Third, after the generated-subsource inheritance principle is stated, the
all-prime boundary closes: it uniquely forces the dense-orbit rule,
source-transitive histories exist, the finite-language criterion is
equivalent to that rule, and every finite packet, hence every finite prime
combination and every mixed fibre, occurs inside each such single history.

Fourth, the arithmetic identification closes only after the topology has
acted.  Atlas-irreducible return types are precisely the original unary
births, so their posterior ranks are precisely \(\mathbb P\).  No divisor
relation, factorization, prime gap, probability law, batch construction, or
coherence filler enters either birth predicate.

The unavoidable new data are explicit rather than hidden.  The unary source
does not itself force permission to revisit an earlier state; that is the one
finite-source enlargement.  Likewise, finite verdicts cannot determine any
infinite-boundary verdict, as Proposition~\ref{prop:finite-not-boundary}
proves, so the generated-subsource inheritance principle is an additional
boundary axiom rather than a theorem secretly extracted from the unary
ledger.  Once those two extensions are accepted, their normal form and
boundary predicate leave no further choice.  There is still no unique
canonical boundary history, and boundary birth cannot be certified from one
finite prefix.  Nevertheless
Theorem~\ref{thm:unique-provenance-class} gives one unique all-prime
source-provenance class; none of these facts leaves a gap in the
prime-spectrum or descent proofs.

\section{Conclusion}

The original one-direction mother supplies the unary primitive types.  The
return-history extension does not glue completed unary births; it enlarges
the source itself so that endpoint returns and their overlap incidence exist
before classification.  Connected finite packets give genuine arbitrary
finite prime births and descend internally to the original unary births.

On the inverse-limit boundary, the generated-subsource inheritance principle
uniquely turns dense shift orbit into the source-native birth condition.  Any
history satisfying it contains every finite born packet as a window, and its
primitive normalized-return spectrum has posterior rank set exactly
\(\mathbb P\).  Repeated occurrences let all finite combinations live inside
the same path without forcing a Boolean face lattice on one copy of each
prime.  The finite theory, the explicit boundary axiom, the all-prime object,
provenance, and the descent tower are then mutually compatible.  Although no
single carrier path is canonical, all all-prime carriers determine the same
unique source-provenance class.
\end{document}
